
\begin{ledgroupsized}[r]{120mm}
\footnotesize
\pstart
\noindent\textbf{Überlieferung:}
\pend
\end{ledgroupsized}
\begin{ledgroupsized}[r]{114mm}
\footnotesize
\pstart
\parindent -6mm
\makebox[6mm][l]{\textit{L}}%
Auszüge mit Bemerkungen aus 
\indexN{\textls{Saccheri}, Girolamo 1667\textendash1733}G.~\textsc{Saccheri},
\cite{00541}\textit{Neo-statica}, Mailand 1708:
LH~XXXV~10, 16~Bl.~3\textendash4. 
Ein Bogen~2\textsuperscript{o};
\href{\Scan{LH_35_10_16_003-004wz}}%
{Wasserzeichen} im Falz. 
Zweieinhalb Seiten auf Bl.~3~r\textsuperscript{o} bis Bl.~4~r\textsuperscript{o};
Bl.~4~v\textsuperscript{o} überliefert N.\;\ref{60124}.
\pend
\end{ledgroupsized}
\selectlanguage{latin}
\vspace{8mm}
\pstart\noindent
\href{\Scan{LH_35_10_16_003r}}%
{\lbrack3~r\textsuperscript{o}\rbrack}
\edtext{\cite{00541}\textit{Neostatica}\indexS{Neostatica}
\textit{auctore Hieronymo Saccherio\indexN{\textls{Saccheri}, Girolamo 1667\textendash1733}
e Soc.~Jesu in Ticinensi universitate Matheseos professore. Mediolani 1708}.}{\lemma{\textit{Neostatica} \lbrack...\rbrack\ \textit{1708}}\Cfootnote{%
\indexN{\textls{Saccheri}, Girolamo 1667\textendash1733}\textsc{G.~Saccheri}, 
\cite{00541}\textit{Neo-statica}, Mailand 1708, Titelseite, mit Auslassungen.}} 
4\textsuperscript{o}. 
\edtext{pagg~168.}{%
\lemma{}%
\Bfootnote{%
pagg~168
\textit{erg.~L}}}
\edtext{In praef. laudat 
\textit{Dn.~Marchionem Pium Bellisomum\indexN{\textls{Bellisomi}, Pio ? \textendash\ nach 1708} 
patricium Ticinensem, inter praeclaras animi dotes mathematicis disciplinis haud leviter instructum}.}{%
\lemma{in praef. \lbrack...\rbrack\ \textit{instructum}}%
\Cfootnote{\cite{00541}a.\,a.\,O., Praefatio, S.~\lbrack VII\rbrack\ mit Auslassungen.}}
\pend 
\pstart
Lib.~1.
\edtext{\textls{impetum}\indexS{impetus} definit a sola velocitate,\indexS{velocitas}}{\lemma{\textls{impetum} \lbrack...\rbrack\ velocitate}%
\Cfootnote{\cite{00541}a.\,a.\,O., lib.~1, def.~9, S.\;2.}}
\edtext{\textls{momentum}\indexS{momentum} a velocitate\indexS{velocitas} et pondere.\indexS{pondus}}{\lemma{\textls{momentum} \lbrack...\rbrack\ pondere}%
\Cfootnote{\cite{00541}a.\,a.\,O., def.~4, S.\;2.}} %
\pend
\begin{edfigure}%
\includegraphics[scale=1]{%
gesamttex/edit_VIII-4/images/LH_35_10_16_003-004_d1_003r.pdf}
\edcaption{\edtext{\lbrack\textit{Fig.~1}\rbrack}{%
\lemma{\lbrack\textit{Fig.~1}\rbrack}%
\Cfootnote[nonum]{\cite{00541}a.\,a.\,O., S.\;3.}}}
\end{edfigure}
\pstart
\edtext{\edtext{Si pondus}{\lemma{Si}\Bfootnote{%
\textit{(1)}~corpus %
\textit{(2)}~pondus~\textit{L}}}
tendat ire linea \textit{AC}, sed cogatur ire linea \textit{AB}, impetum \textit{AC} vocat primigenium,\indexS{impetus primigenius} sed \textit{AB} coactum\indexS{impetus coactus} vel subnascentem.\indexS{impetus subnascens}}{%
\lemma{Si pondus \lbrack...\rbrack\ subnascentem}%
\Cfootnote{\cite{00541}a.\,a.\,O., def.~10, S.\;2\,f.}}
\pend 
\begin{edfigure}%
\includegraphics[scale=1]{%
gesamttex/edit_VIII-4/images/LH_35_10_16_003-004_d2_003r.pdf}
\edcaption{\edtext{\lbrack\textit{Fig.~2}\rbrack}{%
\lemma{\lbrack\textit{Fig.~2}\rbrack}%
\Cfootnote[nonum]{Vorlage ist die Figur \cite{00541}a.\,a.\,O., S.\;4 mit Änderungen.}}}%
\end{edfigure}
\pstart
\edtext{%
\edtext{\textit{Postulat}\lbrack\textit{um}\rbrack.}{%
\lemma{}%
\Bfootnote{%
\textit{Postulat}
\textit{erg.~L, ändert Hrsg.}}}
\textit{Si quoddam pondus \textit{A}, citatum intelligatur impetu secundum \textit{AB}, dum interim} %
\edtext{\textit{ipsum planum}}{\lemma{\textit{ipsum}}\Bfootnote{%
\textit{(1)}~punctum %
\textit{(2)}~\textit{planum}~\textit{L}}}
\textit{movetur ipsi sibi parallelum describente puncto \textit{A}, rectam \textit{AD}, postulamus nullam in pondere sequi} %
\edtext{\textit{variationem quoad}}{\lemma{\textit{variationem}}\Bfootnote{%
\bars{in \textit{streicht Hrsg.}}
\textit{(1)}~quod %
\textit{(2)}~\textit{quoad}
\textit{(3)}~\textit{quoad}~\textit{L}}}
\textit{planum \textit{AB}.}}{%
\lemma{\textit{Postulat}\lbrack\textit{um}\rbrack. \lbrack...\rbrack\ planum \textit{AB}}%
\Cfootnote{\cite{00541}a.\,a.\,O., Postulatum, S.\;4 mit Auslassungen.}}
\edtext{Notat postulatum hoc jam apud alios esse receptum.}{%
\lemma{Notat \lbrack...\rbrack\ receptum}%
\Cfootnote{\cite{00541}a.\,a.\,O., S.\;4.}}
\edtext{Hinc infert prop.~8. si mobile duplici impetu moveatur,\indexS{mobile duplici impetu motum}
secundum latera parallelogrammi\indexS{latus parallelogrammi} moveri in ejus diagonali.\indexS{diagonalis parallelogrammi}}{\lemma{Hinc \lbrack...\rbrack\ diagonali}\Cfootnote{\cite{00541}a.\,a.\,O., prop.~8, S.\;10\,f.}}
\pend
\pstart
Ingeniosus est autor %
\edtext{%
\edtext{lib.~1. prop.~13}{\lemma{}\Bfootnote{lib.~1. prop.~13 \textit{erg.~L}}}
in tentanda demonstratione
\edtext{anguli}{\lemma{}\Bfootnote{anguli \textit{erg.~L}}}
incidentiae et reflexionis.\indexS{angulus incidentiae et reflexionis}}%
{\lemma{lib.~1. \lbrack...\rbrack\ reflectionis}\Cfootnote{\cite{00541}a.\,a.\,O., S.\;15\,f.}} %
\edtext{Tentat etiam in scholio post prop.~13}{\lemma{Tentat \lbrack...\rbrack\ prop.~13}\Cfootnote{\cite{00541}a.\,a.\,O., S.\;17\textendash19.}}
alia via demonstrare 
\edtext{postulatum supra dictum.}{%
\lemma{postulatum supra dictum}%
\Cfootnote{\cite{00541}a.\,a.\,O., S.\;4.}}
Sed illud postulatum praesupponit, mobile\indexS{mobile} a plano cui incumbit abripi; et ita res revera redit ad motum
\edlabel{35_10_16_003-004_2a}%
\edtext{}{%
\xxref%
{35_10_16_003-004_2a}{35_10_16_003-004_2b}%
\lemma{compositum.}%
\Bfootnote{%
\textit{(1)}~Caeterum
\textit{(2)}~Po
\textit{(3)}~Ipse~\textit{L}}}%
compositum.\indexS{motus compositus}
\edtext{Ipse\edlabel{35_10_16_003-004_2b} tamen conatur ex eo deducere, quod non potest intelligi ulla ratio, cur motus mobilis in plano\indexS{motus mobilis in plano} alius sit quam sibi parallelus.}{%
\lemma{Ipse \lbrack...\rbrack\ parallelus}%
\Cfootnote{\cite{00541}a.\,a.\,O., S.\;18.}}
\pend
\begin{edfigure}%
\includegraphics[scale=1]{%
gesamttex/edit_VIII-4/images/LH_35_10_16_003-004_d3_003r.pdf}
\edcaption[0.2em]{\edtext{\lbrack\textit{Fig.~3}\rbrack}{%
\lemma{\lbrack\textit{Fig.~3}\rbrack}%
\Cfootnote[nonum]{\cite{00541}a.\,a.\,O., S.\;19.}}}
\end{edfigure}
\pstart
Aliter etiam ingeniose satis procedit: postulat 
\edtext{%
\textit{si ponderi \textit{A}, citato secundum \textit{AC}, solus relictus sit motus liber\indexS{motus liber}} 
in immobili \textit{BAD}, ferri potius ad partes \textit{B}, quam ad partes 
\edlabel{35_10_16_003-004_3a}%
\edtext{}{%
\xxref%
{35_10_16_003-004_3a}{35_10_16_003-004_3b}%
\lemma{\textit{D} et}\Bfootnote{%
\textit{(1)}~impetum %
\textit{(2)}~fore~\textit{L}}}%
\textit{D}}{%
\lemma{si ponderi \lbrack...\rbrack\ \textit{D}}%
\Cfootnote{\cite{00541}a.\,a.\,O., postul.~1, S.\;19 mit Auslassungen.}}
\edtext{et fore\edlabel{35_10_16_003-004_3b} majorem velocitatem, si angulus \textit{BAC} sit acutior.}{%
\lemma{et fore \lbrack...\rbrack\ acutior}%
\Cfootnote{\cite{00541}a.\,a.\,O., postul.~2, S.\;19.}}
\pend 
\begin{edfigure}%
\includegraphics[scale=1]{%
gesamttex/edit_VIII-4/images/LH_35_10_16_003-004_d4_003r.pdf}
\edcaption{\edtext{\lbrack\textit{Fig.~4}\rbrack}{%
\lemma{\lbrack\textit{Fig.~4}\rbrack}%
\Cfootnote[nonum]{\cite{00541}a.\,a.\,O., S.\;20.}}}
\end{edfigure}
\pstart
\edtext{Item postulat, si mobile tendat directione \textit{AB}, et directione \textit{AC}, ferri directione media angulum secante \textit{AD}.}{%
\lemma{Item \lbrack...\rbrack\ \textit{AD}}%
\Cfootnote{\cite{00541}a.\,a.\,O., postul.~3, S.\;20.}}
\edtext{Postulat etiam motum secundum \textit{AD} fore aequabilem,\indexS{motus aequabilis} si motus secundum \textit{AB}, et secundum \textit{AC} sint aequabiles.}{%
\lemma{Postulat \lbrack...\rbrack\  aequabiles}%
\Cfootnote{\cite{00541}a.\,a.\,O., postul.~4, S.\;20.}}
\pend
\begin{edfigure}%
\includegraphics[scale=1]{%
gesamttex/edit_VIII-4/images/LH_35_10_16_003-004_d5_003r.pdf}
\edcaption{\edtext{\lbrack\textit{Fig.~5}\rbrack}{%
\lemma{\lbrack\textit{Fig.~5}\rbrack}%
\Cfootnote[nonum]{Vorlage ist die Figur \cite{00541}a.\,a.\,O., S.\;21 mit Änderungen.}}}
\end{edfigure}
\pstart
Haereo in
\edtext{demonstratione propositionis 14
\edtext{libri~1. Si}{\lemma{libri~1.}\Bfootnote{%
\textit{(1)}~Im %
\textit{(2)}~Si~\textit{L}}}
pondus\indexS{pondus duobus impetibus citatum} 
citatum sit duobus impetibus\indexS{impetus} \textit{AD} et \textit{AC}, %
\edtext{et impetu \textit{AC}}{\lemma{et}\Bfootnote{%
\textit{(1)}~\textit{AC} %
\textit{(2)}~impetu \textit{AC}~\textit{L}}}
subnascatur impetus\indexS{impetus subnascens} %
\edtext{adaequatus}{\lemma{}\Bfootnote{adaequatus \textit{erg.~L}}} 
\textit{AB} aequalis impetui 
\edtext{\lbrack\textit{AD}\rbrack,}{%
\lemma{}%
\Bfootnote{%
\textit{AB}
\textit{L ändert Hrsg. nach Vorlage}}}
ostendit autor impetum vivum\indexS{impetus vivus} 
(credo hoc nomine intelligit resultantem\indexS{impetus resultans})
non esse in \textit{R} aut (\textit{R}), supponit tamen esse aliquem, 
ergo concludit esse in \textit{AN} normali ad \textit{AB}.%
}{\lemma{demonstratione \lbrack...\rbrack\ ad \textit{AB}}\Cfootnote{a.\,a.\,O., S.\;21\,f.\cite{00541}}}
Sed debebat ostendere esse aliquem, videtur enim esse nullus, si subnascens\indexS{impetus subnascens}
aequali \textit{AD} tollatur, si modo mentem autoris terminos novantem satis intelligo. 
\pend
\pstart
\edtext{In scholio propositionis 26 ad componendos tres impetus\indexS{compositio trium impetuum} 
primum componit\indexS{compositio duorum impetuum} duos simplices\indexS{impetus simplex} 
in unum compositum,\indexS{impetus compositus} et cum hoc tertium 
simplicem,\indexS{impetus simplex}}{\lemma{In scholio \lbrack...\rbrack\ simplicem}\Cfootnote{\cite{00541}a.\,a.\,O., S.\;43\textendash45.}} 
sed
\href{\Scan{LH_35_10_16_003v}}%
{\lbrack3~v\textsuperscript{o}\rbrack}
\edtext{non animadvertit regulam generalem\indexS{regula generalis de compositione motuum} 
quod tunc sit progressus centri gravitatis\indexS{progressus centri gravitatis} per numerum impetuum\indexS{numerus impetuum} multiplicatus.}{%
\lemma{non \lbrack...\rbrack\ multiplicatus}%
\Cfootnote{Leibniz spielt auf seinen Aufsatz 
\cite{00542}\textit{Regle generale de la composition des mouvemens}, \textit{JS} vom 7.~September 1693, S.\;417\textendash419 an, 
dessen Hauptergebnis die angeführte allgemeine Regel ist.
In den entsprechenden Passagen von N.\;\ref{60124} und N.\;\ref{60122} (\refnote{FnC:35_10_16_004_regle} und \refnote{FnC:35_10_16_001-002_regle}) wird die Publikation ausdrücklich zitiert.}}
\pend 
\begin{edfigure}[0.5em]%
\includegraphics[scale=1]{%
gesamttex/edit_VIII-4/images/LH_35_10_16_003-004_d6_003v.pdf}
\edcaption{\edtext{\lbrack\textit{Fig.~6}\rbrack}{%
\lemma{\lbrack\textit{Fig.~6}\rbrack}%
\Cfootnote[nonum]{Vorlage ist die Figur aus 
\indexN{\textls{Saccheri}, Girolamo 1667\textendash1733}\textsc{Saccheri}, 
\cite{00541}\textit{Neo-statica}, S.\;59 mit Änderungen.}}}%
\end{edfigure}
\pstart
\edtext{Si sit 
\edtext{libra\indexS{libra} \textit{CNR},}{%
\lemma{libra \textit{CNR}}%
\Cfootnote{Zu Saccheris Begriff der \glqq libra\grqq\ siehe \cite{00541}a.\,a.\,O., lib.~2, \textit{Definitiones}, S.\;55.}}
et centrum motus\indexS{centrum motus} \textit{N}, %
\edtext{et directiones}{\lemma{et}\Bfootnote{%
\textit{(1)}~mobilium %
\textit{(2)}~directiones~\textit{L}}}
ponderum
\edtext{aequalium\indexS{pondera aequalia}}{\lemma{}\Bfootnote{aequalium \textit{erg.~L}}} 
\textit{C} et \textit{R} %
\edtext{sint in}{\lemma{sint}\Bfootnote{%
\textit{(1)}~versus %
\textit{(2)}~in~\textit{L}}}
\textit{CD} et in \textit{RD}, ait autor lib.~2 %
\edtext{prop.~6. si sint in aequilibrio\indexS{aequilibrium} tunc fore rationem \textit{CN} ad \textit{RN}}%
{\lemma{prop.~6.}\Bfootnote{%
\textit{(1)}~fore %
\textit{(2)}~si sint in aequilibrio tunc fore rationem
\textit{(a)}~ad \textit{RN} %
\textit{(b)}~\textit{CN} ad \textit{RN}~\textit{L}}}
compositam \textit{ex ratione directa \textit{CD}}
\edtext{\textit{ad \textit{RD} et ex reciproca} velocitatum\lbrack,\rbrack\ seu velocitatis}{\lemma{\textit{ad \textit{RD}}}\Bfootnote{%
\textit{(1)}~ex %
\textit{(2)}~\textit{et ex reciproca}
\textit{(a)}~velocitatis %
\textit{(b)}~velocitatum seu velocitatis~\textit{L}}}
\textit{secundum \textit{RD}} ad velocitatem secundum \textit{CD}.}{%
\lemma{Si sit \lbrack...\rbrack\ secundum \textit{CD}}%
\Cfootnote{\cite{00541}a.\,a.\,O., lib.~2, prop.~6, S.\;59 mit Auslassungen.}}
\pend
\pstart
Examinemus hoc ex nostris: %
\edtext{Ex \textit{N} agantur perpendiculares ad \textit{CD}}{\lemma{Ex \textit{N}}\Bfootnote{%
\textit{(1)}~agatur perpendicularis %
\textit{(2)}~agantur perpendiculares ad 
\textit{(a)}~\textit{CN} %
\textit{(b)}~\textit{CD}~\textit{L}}}
et ad \textit{RD}, nempe \textit{NE} et \textit{NF}. Si jam sint in aequilibrio pondera,\indexS{pondera in aequilibrium} erit \textit{NE} in conatum\indexS{conatus}
\edtext{ipsius \textit{C},}{\lemma{ipsius}\Bfootnote{%
\textit{(1)}~\textit{CE}, %
\textit{(2)}~\textit{C},~\textit{L}}}
$=$ ipsi \textit{NF} in conatum ipsius \textit{R}, seu \textit{NE} in conat.~\textit{C} $=NF$ in
\edtext{conat.~\textit{R}. Jam}{\lemma{conat. \textit{R}.}\Bfootnote{%
\textit{(1)}~Ergo %
\textit{(2)}~Jam~\textit{L}}}
triang. \textit{CDN} est ad %
\edtext{triang. \textit{RDN} ut}{\lemma{triang. \textit{RDN}}\Bfootnote{%
\textit{(1)}~ad %
\textit{(2)}~ut~\textit{L}}}
\textit{CN}
\edtext{ad \textit{RN}, seu}{\lemma{ad \textit{RN},}\Bfootnote{%
\textit{(1)}~\bars{vel \textit{streicht Hrsg.}} ut \textit{NF} in \textit{RD}, ut %
\textit{(2)}~seu~\textit{L}}}
ut \textit{NE} in \textit{CD} ad \textit{NF} in \textit{RD}. 
Habemus ergo $CN:RN=NE$ in $CD\mathbin{,:\,,}NF$ in $RD=(NE:NF)\cdot (CD:RD)$. 
Sed $NE:NF= \text{conat.}\; R:\text{conat.}\; C$. 
ex demonstratione mea.
Ergo fit $CN:RN=(\text{conat.}\; R:\text{conat.}\; C)\cdot (CD:RD)$. Et proinde vera est propositio autoris. 
\pend
\pstart
\edtext{%
Hinc infert autor si centrum gravium\indexS{centrum gravium} %
\edtext{sit in \lbrack\textit{N}\rbrack, sitque}{\lemma{sit in}\Bfootnote{%
\bars{\textit{D} \textit{ändert Hrsg.}},
\textit{(1)}~sintque pondera \bars{\textit{C} \textit{R} \textit{erg.}} reciproce ut dista %
\textit{(2)}~sitque~\textit{L}}}
pondus \textit{C} ad pondus \textit{R}, ut \textit{RN} ad \textit{CN}, fore %
\edtext{conatus seu gravitates\indexS{conatus seu gravitas} in ratione}{\lemma{conatus}\Bfootnote{%
\textit{(1)}~in r %
\textit{(2)}~seu %
\textit{(a)}~gravitationes %
\textit{(b)}~gravitatis %
\textit{(c)}~gravitates in ratione~\textit{L}}}
distantiarum\indexS{conatus in ratione distantiarum a centro} a centro.}{%
\lemma{Hinc \lbrack...\rbrack\ centro}%
\Cfootnote{\cite{00541}a.\,a.\,O., prop.~11, S.\;64.}}
(+~in Hypothesi communi ubi 
\edtext{\lbrack centrum\rbrack}{\lemma{}\Bfootnote{%
centrum \textit{gestr.~L, wieder gültig gemacht Hrsg.}}}
concursus\indexS{centrum concursus} seu centrum gravium\indexS{centrum gravium} ponitur abesse
\edtext{infinite, gravitates\indexS{gravitas}}{\lemma{infinite,}\Bfootnote{%
\textit{(1)}~eri %
\textit{(2)}~erit %
\textit{(3)}~etiam %
\textit{(4)}~gravitates~\textit{L}}}
censentur aequales, et distantiae\indexS{distantia a centro} etiam aequales.~\lbrack+)\rbrack
\pend 
\pstart
\edtext{%
\edtext{Postulatum libro 2\textsuperscript{do}}{\lemma{Postulatum}\Bfootnote{%
\textit{(1)}~libri 2\textsuperscript{di} %
\textit{(2)}~libro 2\textsuperscript{do}~\textit{L}}}
tale praemiserat autor: 
\textit{existente \textit{D} centro communi gravium,\indexS{centrum commune gravium} rectus sit}
\edtext{\textit{angulus \textit{RND}}, et data puncta, \textit{R}, \textit{N}, \textit{D}\,\lbrack;\rbrack\ \textit{data}}{\lemma{\textit{angulus \textit{RND}},}\Bfootnote{%
\textit{(1)}~\textit{data} %
\textit{(2)}~\bars{et data puncta, \textit{R}, \textit{N}, \textit{D} \textit{erg.}} \textit{data}~\textit{L}}}
\textit{etiam sint duo quaelibet gravia \textit{R} et \textit{C}, quorum primum constitutum intelligatur in puncto \textit{R},}
alterum in puncto \textit{C}, 
\textit{postulamus posse reperiri in recta \textit{RN} in infinitum} producta, 
\textit{quoddam punctum \textit{C}, in quo si locetur alterum grave \textit{C}, dictorum gravium aequilibrium\indexS{aequilibrium gravium} habeatur in puncto \textit{N}.}}{%
\lemma{Postulatum \lbrack...\rbrack\ \textit{puncto \textit{N}}}%
\Cfootnote{\cite{00541}a.\,a.\,O., S.\;56 mit Auslassungen.}}
\pend
\begin{edfigure}%
\includegraphics[scale=1]{%
gesamttex/edit_VIII-4/images/LH_35_10_16_003-004_d7_003v.pdf}\hspace*{2.2em}
\edcaption{\lbrack\textit{Fig.~7}\rbrack}
\end{edfigure}
\pstart
\edtext{%
Hinc autor contendit probare propositione 15 et ejus corollario fieri non posse, 
ut \textit{impetus naturales gravium\indexS{impetus naturales gravium} ex quiete versus centrum\indexS{impetus gravium versus centrum} commune} gravium,\indexS{centrum commune gravium}
\textit{sint in} quavis \textit{distantia aequales}\lbrack;\rbrack\ 
\textit{et multo minus, quod in majore distantia} (quae tamen mea sententia 
\edtext{est)}{\lemma{}\Bfootnote{est \textit{erg.~L}}}
\textit{sint minores}.%
}{\lemma{Hinc \lbrack...\rbrack\ \textit{minores}}%
\Cfootnote{\cite{00541}a.\,a.\,O., \textit{Corollarium}, S.\;69 mit Auslassungen.}}
\edtext{Sed quid}{\lemma{Sed}\Bfootnote{%
\textit{(1)}~non %
\textit{(2)}~quid~\textit{L}}}
si fingamus \textit{C} trahi filo\indexS{filum} \textit{CDEF}, quod transeat super
\edtext{trochlea\indexS{trochlea} \textit{E}, et}{\lemma{trochlea \textit{E},}\Bfootnote{%
\textit{(1)}~cui %
\textit{(2)}~et~\textit{L}}}
habeat annexum pondus\indexS{pondus filo annexum} \textit{F}, et \textit{R} trahi filo\indexS{filum} \textit{RDGH} %
\edtext{quod transeat super}{\lemma{quod}\Bfootnote{%
\textit{(1)}~transeat %
\textit{(2)}~t %
\textit{(3)}~transeat %
\textit{(a)}~etiam %
\textit{(b)}~super~\textit{L}}}
trochlea\indexS{trochlea} \textit{G}, et habeat annexum pondus\indexS{pondus filo annexum} \textit{H}. %
\pend
\pstart
Data sunt %
\edtext{pondera \textit{F}}{\lemma{pondera}\Bfootnote{%
\textit{(1)}~\textit{H} pro \textit{R}, et %
\textit{(2)}~\textit{F}~\textit{L}}}
pro \textit{C}, et \textit{H} pro \textit{R}, ponamus praeterea datas esse eorum velocitates, vel gravitationes %
\edtext{qualescunque, sive ponamus dari}{\lemma{qualescunque,}\Bfootnote{%
\textit{(1)}~vel d %
\textit{(2)}~sive ponamus %
\textit{(a)}~sup %
\textit{(b)}~dari~\textit{L}}}
per celeritates vel per plana inclinata\indexS{planum inclinatum} \textit{I} et \textit{K}. Dico fore, \textit{N} 
\edtext{punctum aequilibrii,\indexS{punctum aequilibrii}}{%
\lemma{punctum}%
\Bfootnote{%
\textit{(1)}~aequil.
\textit{(2)}~aequilibrii,~\textit{L}}}
si datis %
\edtext{punctis \textit{R} et \textit{N},}{\lemma{punctis}\Bfootnote{%
\textit{(1)}~\textit{N} %
\textit{(2)}~\textit{R} et \textit{N},~\textit{L}}}
itemque puncto \textit{D} ita ut angulus \textit{RND} sit rectus\lbrack,\rbrack\ %
\edtext{adeoque dato et puncto \textit{X},}{\lemma{}\Bfootnote{adeoque dato et puncto \textit{X}, \textit{erg.~L}}} 
assumatur punctum \textit{C} tale, ut sit \textit{NL} ad \textit{NX} in ratione momenti\indexS{momentum} ipsius \textit{H} ad momentum ipsius \textit{F}. Patet hic %
\edtext{dato centro}{\lemma{dato}\Bfootnote{%
\textit{(1)}~puncto %
\textit{(2)}~centro~\textit{L}}}
gravium,\indexS{centrum gravium} %
\edtext{posse rationem momenti\indexS{momentum} ponderis appensi ipsi \textit{R} datae extremitati librae ad alterius}{\lemma{posse}\Bfootnote{%
\textit{(1)}~momentum ponderis appensi ad extremitatem librae 
\textit{(a)}~datae
\textit{(b)}~datam \textit{R}, assumi utcunque %
\textit{(2)}~rationem momenti ponderis appensi \bars{ipsi \textit{R} \textit{erg.}}
datae extremitati librae \bars{\textit{R} \textit{streicht Hrsg.}} ad alterius~\textit{L}}}
momentum\indexS{momentum} assumi utcunque parvam, quia data \textit{NX}, potest \textit{NL} fieri utcunque parva, si satis
\href{\Scan{LH_35_10_16_004r}}%
{\lbrack4~r\textsuperscript{o}\rbrack} 
admoveatur \textit{C} ad \textit{N}. Sed non potest assumi ratio utcunque magna, quia non potest esse major
\edtext{quam \textit{ND}}{\lemma{quam}\Bfootnote{%
\textit{(1)}~\textit{NL} %
\textit{(2)}~\textit{ND}~\textit{L}}}
ad \textit{NX}\lbrack,\rbrack\ utcunque enim removeas \textit{C} ab \textit{N}, nunquam \textit{NL} excedet \textit{ND}. 
\pend
\pstart
Ponamus pondera \textit{H} et \textit{F} aequalia,\indexS{pondera aequalia} et sint %
\edtext{momenta,\indexS{momentum} seu}{\lemma{momenta}\Bfootnote{%
\textit{(1)}~recipro %
\textit{(2)}~seu %
\textit{(a)}~velocit %
\textit{(b)}~ve %
\textit{(5)}~, seu~\textit{L}}}
hoc loco conatus\indexS{conatus} reciproce ut distantiae, %
\edtext{seu sit conatus}{\lemma{seu sit}\Bfootnote{%
\textit{(1)}~momentum %
\textit{(2)}~conatus~\textit{L}}}
\edtext{applicatus ipsi \textit{R}, ad conatum applicatum ipsi \textit{C}, ut \textit{CD}}{\lemma{applicatus}\Bfootnote{%
\textit{(1)}~ad \textit{R}, ut conatus applicatus ad \textit{C} %
\textit{(2)}~ipsi \textit{R}, ad conatum applicatum ipsi \textit{C}, %
\textit{(a)}~ut %
\textit{(aa)}~\textit{DC} ad %
\textit{(bb)}~\textit{CD} ad %
\textit{(b)}~ut \textit{CD}~\textit{L}}}
ad \textit{RC} sive $\text{con.}\;R:\text{con.}\;C=CD:RD$. Ita suppono. 
Rursus $NL:NX=\text{con.}\;R:$ %
\edtext{$\text{con.}\;C$. Ex}{\lemma{con.~\textit{C}.}\Bfootnote{%
\textit{(1)}~Fit %
\textit{(2)}~Ex~\textit{L}}}
\edtext{meis demonstratis.}{%
\lemma{meis demonstratis}%
\Cfootnote{%
Siehe bspw. die \textit{Solutio problematis R.\,P.\;Vannii S.\,J.} (April bis Juni 1690), N.\;\ref{41823} in diesem Band, bes. \refpassage{35_10_13_010-011_17a}{35_10_13_010-011_17b}.}}
Ergo fit
$CD:RD=NL:NX$\lbrack,\rbrack\ sed per priora $NL:NX=(CN:RN)\cdot (RD:CD)$. Ergo fiet \edlabel{35_10_16_003-004_1a}$CN:RN=\text{qu.}\;CD:\text{qu.}\;RD$.%
\edtext{}{{\xxref{35_10_16_003-004_1a}{35_10_16_003-004_1b}}\lemma{qu.~\textit{RD}.}\Bfootnote{%
\textit{(1)}~$NR\,\text{qu}+ND\,\text{qu}=$ %
\textit{(2)}~\bars{$CD\,\text{qu}$ \textit{streicht Hrsg.}}
\textit{(3)}~$CN:RN=CN\,\text{qu}+ND\,\text{qu}$~\textit{L}}}
\pend
\pstart
$CN:RN=CN\,\text{qu}+ND\,\text{qu}\mathbin{,:\,,}RN\,\text{qu}+ND\,\text{qu}$\edlabel{35_10_16_003-004_1b}. 
Fiet $CN\cdot RN\,\text{qu}+CN\cdot ND\,\text{qu}=RN\cdot CN\,\text{qu}+RN\cdot ND\,\text{qu}$. 
\edlabel{35_10_16_003-004_5a}%
\edtext{}{%
\xxref%
{35_10_16_003-004_5a}{35_10_16_003-004_5b}%
\lemma{Quaeritur \textit{CN}.}%
\Bfootnote{%
\bars{Huic \lbrack...\rbrack\ $+ND\,\text{qu}$. \textit{erg.}}
\textit{(1)}~$CN\,\text{qu}-\displaystyle\frac{RN\,\text{qu}+ND\,\text{qu}}{RN}%
\rule[-2.4ex]{0pt}{6.2ex}%
+\displaystyle\efrac{\text{qu}(RN\,\text{qu}+}{}$ %
\textit{(2)}~$CN\,\text{qu}-$~\textit{L}}}%
Quaeritur \textit{CN}. Huic aequationi satisfit ponendo $RN=CN$. 
Fit enim $RN\,\text{qu}+ND\,\text{qu}=RN\,\text{qu}+ND\,\text{qu}$.
\pend
\pstart\noindent
\rule[-2.8ex]{0pt}{7ex}%
$CN\,\text{qu}-\edlabel{35_10_16_003-004_5b}%
\displaystyle\frac{RD\,\text{qu}}{RN}CN+\displaystyle\frac{RD.\text{biqu}}{4RN\,\text{qu}}=\displaystyle\frac{RD\,\text{biqu}}{4RN\,\text{qu}}-ND\,\text{qu}$.
\pend
\pstart\noindent
\rule[-2.9ex]{0pt}{6.5ex}%
$CN=\displaystyle\frac{RD\,\text{qu}+\surd (RD\,\text{biqu}-4\,\text{qu}\,(RD\cdot RN)\lbrack)\rbrack}{2RN}$. 
\pend
\pstart
Ergo res procedit %
\edtext{si $RD\,\text{qu}$ non sit minus quam}{\lemma{si $RD\,\text{qu}$}\Bfootnote{%
\textit{(1)}~sit majus quam %
\textit{(2)}~non sit minus quam~\textit{L}}}
bis~$RN\cdot ND$, sed hoc semper verum est %
\edtext{nam nunquam}{\lemma{nam}\Bfootnote{%
\textit{(1)}~semper %
\textit{(2)}~nunquam~\textit{L}}}
quadratum %
\edtext{Hypotenusae minus est}{\lemma{Hypotenusae}\Bfootnote{%
\textit{(1)}~majus est %
\textit{(2)}~minus est~\textit{L}}}
duplo triangulo rectangulo. 
\pend
\pstart
\edlabel{35_10_16_003-004_4a}%
\edtext{}{%
\xxref%
{35_10_16_003-004_4a}{35_10_16_003-004_4b}%
\lemma{Sed}\Bfootnote{%
\textit{(1)}~nunquam stabiliamus %
\textit{(2)}~nunc adhibeamus hypothesin veram, 
\textit{(a)}~seu
\textit{(b)}~mea sententia, seu ut 
\textit{(aa)}~quadrata sin %
\textit{(bb)}~co %
\textit{(cc)}~attractiones~\textit{L}}}%
Sed nunc adhibeamus 
\edtext{hypothesin veram, mea sententia, seu ut attractiones\indexS{attractiones reciprocae quadratis distantiarum}\edlabel{35_10_16_003-004_4b} sint reciproce ut quadrata distantiarum.\indexS{quadratum distantiae}}{%
\lemma{hypothesin \lbrack...\rbrack\ distantiarum}%
\Cfootnote{%
Zu dieser These siehe die unveröffentlichte sog. \glqq zweite Bearbeitung\grqq\ 
des \cite{00543}\textit{Tentamen de motuum coelestium causis}, bspw. LH~XXXV~9, 2~Bl.~54~v\textsuperscript{o} und \cite{01043}\textit{LMG}~VI, S.\;165\,f.
Vgl. dazu auch die Anmerkung zu N.\;\ref{60124}, \refnote{FnC:35_10_16_004_Tentamen}.}}
Fiat $\text{con.}\;R:\text{con.}\;C$ 
\edtext{\lbrack$=$\rbrack}{%
\lemma{}%
\Bfootnote{%
$=$ 
\textit{erg. Hrsg.}}}
$\text{qu.}\;CD:\text{qu.}\;RD=NL:NX=\displaystyle\frac{CN}{RN}\cdot \displaystyle\frac{RD}{CD}$.%
\rule[-2.8ex]{0pt}{6.5ex}
Fiet $CN:RN=CD^3:RD^3$. 
\pend 
